
\chapter[APÊNDICE B -- CRONOGRAMA DE DESENVOLVIMENTO]{{\centering APÊNDICE B \textendash{} CRONOGRAMA DE DESENVOLVIMENTO}}  \label{ap:cronograma}

Este apêndice apresenta a distribuição temporal das atividades do trabalho, organizadas nas macroetapas descritas na Seção~\ref{sec:etapas-pesquisa}, ao longo dos doze meses que compreendem o TCC~I e o TCC~II. A Tabela~\ref{tab:cronograma} preserva as marcações temporais da pré-proposta e ajusta a descrição da avaliação ao escopo atual. A tabela registra o planejamento e não comprova a execução de cada atividade no mês indicado. Cada coluna numerada corresponde a um mês corrido das duas etapas, e cada célula marcada com ``X'' indica a previsão de execução da atividade no período.

\begin{table}[htbp]
    \centering
    \caption{Cronograma de desenvolvimento do trabalho de conclusão de curso.}
    \label{tab:cronograma}
    \footnotesize
    \setlength{\tabcolsep}{4pt}
    \renewcommand{\arraystretch}{1.2}
    \resizebox{\textwidth}{!}{%
    \begin{tabular}{|p{6.5cm}|c|c|c|c|c|c|c|c|c|c|c|c|}
        \hline
        \multirow{2}{*}{\textbf{Atividade}} & \multicolumn{6}{c|}{\textbf{TCC I}} & \multicolumn{6}{c|}{\textbf{TCC II}} \\
        \cline{2-13}
         & \textbf{01} & \textbf{02} & \textbf{03} & \textbf{04} & \textbf{05} & \textbf{06} & \textbf{07} & \textbf{08} & \textbf{09} & \textbf{10} & \textbf{11} & \textbf{12} \\
        \hline
        Revisão Bibliográfica & X & X & X & X & & & & & & & & \\
        \hline
        Escrita da Introdução e Revisão de Literatura & & X & X & X & & & & & & & & \\
        \hline
        Predefinição de escopo e Especificação dos comandos dinâmicos & & & X & X & & & & & & & & \\
        \hline
        Documentação dos procedimentos de desenvolvimento & & & X & & X & X & & & & & & \\
        \hline
        Arquitetura e Modelagem da Plataforma & & & X & X & & & & & & & & \\
        \hline
        Customização dinâmica de \textit{tokens} e sintaxe & & & & X & X & & & & & & & \\
        \hline
        Execução da Análise Léxica e Sintática com o código da IDE & & & & & X & X & & & & & & \\
        \hline
        Execução do código intermediário no Terminal da IDE com \textit{feedback} (avisos e erros) & & & & & & X & X & X & & & & \\
        \hline
        Finalização de fluxos, acessibilidade e responsividade & & & & & & X & X & X & & & & \\
        \hline
        Testes técnicos, preparação e aplicação do protocolo de avaliação & & & & & & & & & X & X & X & \\
        \hline
        Escrita do TCC (Resultados, Discussão e Conclusão) & & & & & & & & & X & X & X & X \\
        \hline
    \end{tabular}%
    }
    \legend{Fonte: elaborado pelo autor.}
\end{table}
